\chapter{Introduction}
Examinations in Indian engineering education follow strong, repeatable patterns. University mid-semester and
end-semester papers draw on a fixed syllabus with a familiar section structure, and the Graduate Aptitude Test in
Engineering (GATE)---taken by several lakh candidates every year for postgraduate admission and public-sector
recruitment---has used a stable 65-question, 100-mark format for more than a decade, mixing multiple choice (MCQ),
multiple select (MSQ) and numerical answer type (NAT) questions. Because the pattern is stable, previous year questions
(PYQs) are the single most used preparation resource: students solve them to learn what is asked, how often, and in
what form.

Using PYQs well is nevertheless laborious. A student who wants to know which subjects of the GATE Computer Science
syllabus carry the most marks, which topics have appeared every year, or whether numerical questions are becoming more
common must read and tally hundreds of questions across many papers. Once the pattern is understood, a second problem
appears: the supply of fresh practice papers that follow the same pattern is limited, and re-solving the same PYQs
quickly loses value because answers are remembered.

Natural Language Processing (NLP) offers tools for the first problem---segmenting documents, representing questions as
vectors, classifying them into syllabus topics and cognitive levels, and measuring similarity across years. Generative
Artificial Intelligence (GenAI), in particular large language models (LLMs) such as Google Gemini, offers a route to the
second: writing new questions on demand. This project, \textbf{PaperMind}, combines the two. It analyses a corpus of past
papers to learn the examination pattern and then uses that learned pattern to steer an LLM that writes new papers,
validated and checked, for whichever examination and subject a student asks for. A conversational study assistant ties
both together and personalises preparation using each student's measured performance.

\section{Challenges}
\begin{itemize}
  \item \textbf{Heterogeneous source documents.} Past papers arrive as digital PDFs, scanned images and photographs. Their
        layouts changed over the years (offline booklets with four codes before 2014, online computer-based tests later),
        equations are often destroyed by PDF text extraction, and several papers bundle more than one session.
  \item \textbf{No labelled data.} Questions are not annotated with syllabus topics or cognitive levels, so supervised
        classification has no ready training set.
  \item \textbf{Faithfulness to the pattern.} A useful generated paper must match the real distribution of subjects,
        marks and question types, not merely be on-topic.
  \item \textbf{Correctness of generated answers.} LLMs produce fluent but occasionally wrong questions and answer keys; a
        practice paper with a wrong key actively misleads a student.
  \item \textbf{Novelty.} Generated questions must not simply reproduce past questions, yet they should be grounded in
        them.
  \item \textbf{Cost and deployment.} Free LLM tiers allow only a small number of requests per day, and free hosting
        offers little memory and CPU, which rules out heavy deep-learning models at run time.
\end{itemize}

\section{Motivation}
Coaching material that summarises ``important topics'' is usually compiled by hand, is rarely transparent about how it
was derived, and is not tailored to an individual. At the same time, general-purpose chatbots will happily produce
``GATE-style questions'' but have no grounding in what GATE actually asks and no check on their answers. A system that
derives the pattern from the papers themselves, explains it with numbers, generates practice material constrained by
that pattern, and verifies what it generates would make PYQ-based preparation faster, more objective and available to
any student with a browser.

\section{Problem Statement}
Given a collection of past examination papers in arbitrary formats, design and implement a system that (i) extracts and
structures every question, (ii) labels each question with its syllabus subject, topic and Bloom's taxonomy level,
(iii) derives the historical pattern of the examination---subject weightage, topic importance, trends, question-type
mix and recurring concepts---and (iv) generates, from a student's natural-language request, a new question paper or
practice set that follows this pattern, carries correct answer keys and is novel with respect to the past papers;
and finally (v) provides a personal assistant that uses these capabilities to guide an individual student's preparation.

\section{Objectives}
\begin{enumerate}
  \item To build a robust ingestion pipeline for digital PDFs, scanned PDFs and page photographs, including joining GATE
        papers with their official answer keys.
  \item To assemble a corpus of GATE CS (2011--2026) and GATE DA (2024--2026) papers with documented provenance.
  \item To label questions with syllabus topics without manual annotation, by distilling LLM labels into a lightweight
        classifier, and to evaluate it on a held-out year.
  \item To compute interpretable historical statistics: weightage, recency-weighted hot topics, trends and recurrence.
  \item To generate pattern-faithful papers with an LLM using retrieval of past questions as exemplars, and to validate
        format, marks, novelty and answer-key correctness.
  \item To build a tool-using conversational study assistant with per-student progress tracking.
  \item To deploy the complete system as a web application at zero hosting cost.
\end{enumerate}

\section{Contributions}
\begin{itemize}
  \item A documented, cleaned corpus of 31 papers (1,971 questions) with official keys where available, including
        detection and correction of mislabelled source files.
  \item A teacher--student labelling scheme in which an LLM labels the corpus and a TF-IDF + logistic-regression model
        learns from it, reaching 77.7\% topic and 88.5\% subject accuracy on unseen 2026 GATE CS papers.
  \item A history-weighted generation pipeline combining blueprint inference, importance-proportional topic sampling,
        retrieval-augmented prompting and blind re-solving of answer keys.
  \item A personal assistant that grounds its advice in corpus statistics and prioritises topics by
        exam weight $\times$ measured weakness.
\end{itemize}

\section{Organisation of the Report}
Chapter~\ref{ch:lit} reviews related work on question analysis, question generation and LLM-based systems.
Chapter~\ref{ch:system} describes the system model: data collection, ingestion, the NLP analysis layer, the generation
pipeline, the assistant and deployment. Chapter~\ref{ch:results} presents the implementation, the corpus analysis, the
classifier evaluation and the generation results. Chapter~\ref{ch:conc} concludes and outlines future work. The appendix
collects the complete analysis reports and sample outputs of the system.

\section{Chapter Summary}
This chapter introduced the role of previous year papers in examination preparation, the twin problems of pattern
discovery and practice-paper scarcity, and the challenges that make automation non-trivial. It stated the problem,
the objectives and the contributions of PaperMind.
